\documentclass[11pt]{article}

\usepackage{amsmath,amssymb,amsthm}
\usepackage[ruled,vlined,linesnumbered]{algorithm2e}
\usepackage{hyperref}
\usepackage{booktabs}
% \usepackage{natbib}  % unused — no bibliography in this draft
\usepackage{fullpage}
\usepackage{microtype}
\usepackage{bm}

\newtheorem{theorem}{Theorem}[section]
\newtheorem{lemma}[theorem]{Lemma}
\newtheorem{corollary}[theorem]{Corollary}
\newtheorem{proposition}[theorem]{Proposition}
\newtheorem{definition}[theorem]{Definition}
\newtheorem{remark}[theorem]{Remark}

\DeclareMathOperator{\poly}{poly}
\DeclareMathOperator{\erfc}{erfc}
\DeclareMathOperator{\corr}{corr}
\newcommand{\bigO}{\mathcal{O}}
\newcommand{\ip}[2]{\langle #1,\, #2 \rangle}

\title{Fast Tournament Equity Computation via \\
       Generating-Function Quadrature and FFT-Accelerated Subproduct Trees}

\author{
  Sam Rosenstrauch \\
  \texttt{sarose550@gmail.com}
}

\date{\today}

\begin{document}
\maketitle

\begin{abstract}
Multi-table poker tournaments present a deceptively hard computational
problem: given each player's chip stack and a prize structure, what is
each player's expected winnings? The standard answer--the Independent
Chip Model (ICM)--requires, in principle, considering every possible
order in which players could be eliminated. For $n$ players this means
$n!$ orderings, which becomes astronomically large even for modest field
sizes. The best known exact algorithm reduces this to $\mathcal{O}(n^2
2^n)$ time via dynamic programming, but remains limited to roughly
$n \leq 23$ players in practice. Approximate methods based on random
sampling exist for larger fields, but converge slowly and cannot achieve
high precision.

I present a deterministic algorithm that sidesteps this combinatorial
explosion entirely. The key insight is that each player's expected payout
can be written as a one-dimensional integral of the coefficients of a generating function that factors as a product over all other players. This integral is
approximated via Gaussian quadrature with $Q = 256$ nodes, achieving
double-precision accuracy ($<5 \times 10^{-12}$ relative error). The
central computational challenge is evaluating this product efficiently
for all $n$ players simultaneously; I address this with an
Fast Fourier Transform (FFT) -accelerated binary subproduct tree, reducing the per-quadrature-point
cost from $\bigO(nk)$ to $\bigO(n \log^2 k)$.

The result is an algorithm that computes all $n$ player equities in
$\mathcal{O}(Qn\log^2 k)$ time, where $k$ is the number of payout
positions--a super-exponential speedup over exact dynamic programming.
I provide implementations tuned to the Apple M3~Max (ARM64) and the AMD
Ryzen~9 7950X (Zen~4 with AVX-512), with a roadmap for porting to the
NVIDIA H200 GPU. On both platforms, all 8192 player equities with
$k = 8192$ payout positions are computed in under 350\,ms
single-threaded and under 37\,ms with 16-thread parallelism.
\end{abstract}

\section{Introduction}\label{sec:intro}

Multi-table poker tournaments are structured competitions in which
players begin with equal chip stacks and are progressively eliminated as
they lose their chips. A payout structure $(\pi_1, \pi_2, \ldots,
\pi_k)$ specifies the prize awarded to the player finishing in each
position, with $k \leq n$ positions receiving nonzero prizes. The
\emph{Independent Chip Model} (ICM) provides a principled mapping from a
player's chip stack to their expected tournament equity--the expected
prize money they will receive.

Despite its central role in tournament strategy, exact ICM computation
is fundamentally difficult because the number of possible elimination
orderings grows factorially with the number of players. Consider even a
modest 10-player table: there are $10! = 3{,}628{,}800$ possible
orderings to sum over. At 20 players this reaches $2.4 \times 10^{18}$.
Bitmask dynamic programming reduces this to $\bigO(n \cdot 2^n)$ time
and $\bigO(2^n)$ space by iterating over subsets of surviving players,
but remains exponential and is limited to roughly $n \leq 23$ in
practice. Monte Carlo sampling algorithms provide approximate solutions for larger fields, but their error shrinks only as $\bigO(1/\sqrt{N_{\mathrm{samples}}})$, meaning that halving the error requires quadrupling the number of samples. High-precision results are effectively out of reach.

The key observation that enables a better approach is that the ICM
equity of each player can be expressed as a one-dimensional integral  rather than an exponential sum. This follows from a classical probabilistic trick: the sequential elimination process (each player's probability of being eliminated next is proportional to their chip stack) is equivalent to assigning each player an independent exponential random variable representing their ``elimination clock,'' and eliminating players
in the order their clocks fire. Under this \emph{exponential clock framing}, the combinatorial sum over orderings collapses into an integral that depends only on the chip stacks.

This integral is then evaluated by Gaussian quadrature--a numerical
integration technique that approximates a definite integral by evaluating
the integrand at $Q$ specially chosen points and combining the results
with precomputed weights. For smooth integrands (as ours is, after an
appropriate change of variables), Gaussian quadrature achieves
exponential convergence in $Q$: with $Q = 256$ points we obtain
double-precision accuracy ($<5 \times 10^{-12}$ relative error), and
adding more points yields diminishing returns. The remaining challenge
is evaluating the integrand efficiently at each quadrature point, which
requires computing a degree-$k$ polynomial that is a product of $n-1$
linear factors--once for each of the $n$ players.

A na\"{i}ve approach would compute each player's polynomial independently
at a cost of $\bigO(nk)$ per quadrature point. I instead use a
\emph{binary subproduct tree} to share computation across all $n$
players simultaneously, reducing the per-point cost to
$\bigO(n \log^2 k)$. The tree structure is analogous to a tournament
bracket: rather than computing each player's product from scratch,
intermediate products are computed once and reused. Each polynomial
multiplication within the tree is performed using the Fast Fourier
Transform (FFT), which reduces the cost of multiplying two polynomials
of degree $k$ from $\bigO(k^2)$ to $\bigO(k \log k)$ by working in the
frequency domain.

My contributions are:
\begin{enumerate}
  \item A rigorous derivation of ICM equities as coefficient extractions
    from a generating function evaluated via numerical quadrature
    (Section~\ref{sec:algorithm}).
  \item A hybrid engine combining a subproduct tree for inter-block
    computation with direct polynomial arithmetic for intra-block
    computation, with cost-model-driven dispatch
    (Section~\ref{sec:algorithm}).
  \item Detailed platform-specific optimizations for ARM64 (Apple M3~Max)
    and x86-64 (AMD Zen~4) microarchitectures, including vDSP/AMX dispatch
    and FFTW/MKL dual-library selection (Section~\ref{sec:experiments}).
  \item A roadmap for GPU acceleration on the NVIDIA H200
    (Section~\ref{sec:gpu}).
\end{enumerate}
\section{Preliminaries}\label{sec:prelim}
% ======================================================================

\subsection{Tournament Model}

\begin{definition}[Tournament instance]
A \emph{tournament instance} is a tuple $(n, \bm{S}, \bm{\pi})$ where
$n \in \mathbb{N}$ is the number of players, $\bm{S} = (S_1, \ldots, S_n)
\in \mathbb{R}_{>0}^n$ is the vector of chip stacks with $S_i > 0$ for all
$i$, and $\bm{\pi} = (\pi_1, \ldots, \pi_k) \in \mathbb{R}_{\geq 0}^k$
is the payout vector with $k \leq n$ and $\pi_1 \geq \pi_2 \geq \cdots
\geq \pi_k \geq 0$.
\end{definition}

\begin{definition}[Malmuth-Harville model]\label{def:mh}
Under the Malmuth-Harville (MH) model, the probability that player $i$ is
the first to be eliminated from a set $T \subseteq \{1, \ldots, n\}$ of
surviving players is:
\begin{equation}\label{eq:mh-elim}
  \Pr[i \text{ elim.\ first} \mid T] = \frac{S_i}{\sum_{j \in T} S_j}\,.
\end{equation}
Elimination proceeds sequentially: after the first player is eliminated, the
model is applied recursively to the remaining set $T \setminus \{i\}$ with
unchanged stacks.
\end{definition}

\begin{definition}[ICM equity]
The \emph{ICM equity} of player $i$ in a tournament $(n, \bm{S},
\bm{\pi})$ is their expected payout under the MH model:
\begin{equation}\label{eq:icm-equity}
  E_i = \sum_{m=1}^{k} \pi_m \cdot \Pr[i \text{ finishes in position } m]\,.
\end{equation}
\end{definition}

\subsection{Prior Algorithms and the Exponential Clock Framing}

The na\"ive algorithm for exact ICM enumerates all $n!$ elimination
orderings and sums each player's payout weighted by the ordering's
probability under the MH model. Bitmask dynamic programming improves this
to $\bigO(n \cdot 2^n)$ time by computing, for each subset $T$ of
surviving players, the probability of reaching that subset and the
conditional placement probabilities. This is the fastest known exact
method, but the exponential state space limits it to roughly $n \leq 23$.
For larger fields, the standard approach is Monte Carlo simulation: sample
random elimination orderings by repeatedly selecting the next player to be
eliminated with probability proportional to their stack
(Equation~\ref{eq:mh-elim}), record each player's finishing position, and
average the resulting payouts. This requires $\bigO(n)$ work per sample and
$\bigO(N_{\mathrm{samples}})$ samples for
$\bigO(1/\sqrt{N_{\mathrm{samples}}})$ relative error--prohibitively slow
for high-precision results.

The key observation that enables a deterministic algorithm is that the MH
model has an equivalent \emph{continuous-time} formulation. Assign each
player~$j$ an independent exponential random variable $T_j \sim
\mathrm{Exp}(S_j)$ representing their elimination time. Players are
eliminated in order of increasing $T_j$. The memoryless property of the
exponential distribution ensures that the conditional probability of any
surviving player being eliminated next is exactly proportional to their
stack, recovering Equation~\ref{eq:mh-elim}. Under this framing, the
probability that player~$i$ finishes in position $r$ is:
\begin{equation}\label{eq:position-prob}
  \Pr[i \text{ finishes in position } r]
  = \int_0^\infty S_i \, e^{-S_i t}
    \!\!\!\sum_{\substack{A \subseteq [n] \setminus \{i\} \\ |A| = r-1}}
    \prod_{j \in A} e^{-S_j t}
    \prod_{j \notin A,\, j \neq i} \!\bigl(1 - e^{-S_j t}\bigr)\, dt\,,
\end{equation}
where $A$ is the set of $r-1$ players whose clocks have \emph{not} fired
by time~$t$ (they survive past player~$i$, finishing in positions $1$
through $r-1$), and the remaining $n-r$ players have been eliminated
earlier.

\subsection{Generating-Function Formulation}

Equation~\ref{eq:position-prob} can be rewritten compactly using generating
functions. Substituting $v = e^{-t}$ (so $dt = -dv/v$ and
$e^{-S_j t} = v^{S_j}$), define:

\begin{definition}[Player generating function]
For $v \in (0, 1)$, let $a_j(v) = v^{S_j}$ and
$b_j(v) = 1 - v^{S_j}$. The \emph{leave-one-out generating function} for
player $i$ is:
\begin{equation}\label{eq:gen-func}
  Q_i(x; v) = \prod_{j \neq i} \bigl(a_j(v) + b_j(v) \, x\bigr)\,.
\end{equation}
The coefficient $[x^m]\, Q_i(x; v)$ is the sum over all subsets of $m$
other players of the product of their elimination probabilities $b_j(v)$
times the remaining players' survival probabilities $a_j(v)$--exactly the
inner sum in Equation~\ref{eq:position-prob} with $m = n - r$ eliminated
players.
\end{definition}

Under the change of variables $v = e^{-t}$, Equation~\ref{eq:position-prob}
becomes:
\begin{equation}\label{eq:position-gf}
  \Pr[i \text{ finishes in position } r]
  = \int_0^1 S_i \, v^{S_i - 1}\, [x^{n-r}]\, Q_i(x; v) \, dv\,.
\end{equation}
Multiplying by $\pi_r$ and summing over $r = 1, \ldots, k$:
\begin{equation}\label{eq:equity-integral}
  E_i = \int_0^1 S_i \, v^{S_i - 1}
    \sum_{r=1}^{k} \pi_r \cdot [x^{n-r}]\, Q_i(x; v) \, dv\,.
\end{equation}
I define the per-quadrature-point inner product:
\begin{equation}\label{eq:inner-product}
  I_i(v) = \sum_{r=1}^{k} \pi_r \cdot [x^{n-r}]\, Q_i(x; v)\,.
\end{equation}

The integral is approximated by the trapezoidal rule after a change of
variables $v = \Phi(y)$ (where $\Phi$ is the standard normal CDF), which
maps $(0,1)$ to $\mathbb{R}$ and yields a rapidly decaying integrand:
\begin{equation}\label{eq:quadrature}
  E_i \approx \sum_{q=1}^{Q} w_q \cdot S_i \cdot a_i(v_q) \cdot v_q^{-1}
    \cdot I_i(v_q)\,,
\end{equation}
where $(v_q, w_q)_{q=1}^Q$ are the quadrature nodes and weights.

\subsection{Polynomial Operations and the Convolution--Correlation Adjoint}
\label{sec:conv-adj}

We work throughout with \emph{truncated polynomials} in the vector space
$\mathbb{R}[x]_{<k} \cong \mathbb{R}^k$ of polynomials of degree
$< k$, equipped with the Euclidean inner product
\begin{equation}\label{eq:poly-ip}
  \ip{f}{g} = \sum_{m=0}^{k-1} f_m \, g_m\,.
\end{equation}

\begin{definition}[Truncated polynomial multiplication (convolution)]
For $f, h \in \mathbb{R}[x]_{<k}$, define the truncated convolution
$\mu_h(f) = (f * h) \bmod x^k$ by
\begin{equation}\label{eq:trunc-conv}
  \bigl(\mu_h(f)\bigr)_m
  = \sum_{\substack{i+j=m \\ 0 \le i,j < k}} f_i \, h_j\,,
  \qquad m = 0, \ldots, k-1.
\end{equation}
The map $\mu_h \colon \mathbb{R}^k \to \mathbb{R}^k$ is linear in~$f$.
\end{definition}

The central algebraic fact underlying the propagation phase is that the
adjoint of convolution is cross-correlation.

\begin{definition}[Cross-correlation]
For $g, h \in \mathbb{R}[x]_{<k}$, the cross-correlation
$\mu_h^*(g) = \corr(g, h)$ is defined by
\begin{equation}\label{eq:corr-def}
  \bigl(\corr(g, h)\bigr)_j
  = \sum_{m=0}^{k-1-j} h_m \, g_{m+j}\,,
  \qquad j = 0, \ldots, k-1.
\end{equation}
Equivalently, $\corr(g, h)$ is the polynomial whose $j$-th coefficient is
the inner product of $h$ with the $j$-shifted truncation of~$g$.
\end{definition}

\begin{lemma}[Correlation is the adjoint of convolution]
\label{lem:adj-conv}
For all $f, g, h \in \mathbb{R}[x]_{<k}$,
\begin{equation}\label{eq:adj-identity}
  \ip{\mu_h(f)}{g} = \ip{f}{\mu_h^*(g)}\,,
  \qquad\text{i.e.,}\qquad
  \ip{f * h}{g} = \ip{f}{\corr(g, h)}\,.
\end{equation}
\end{lemma}
\begin{proof}
Expanding the left-hand side using~\eqref{eq:poly-ip}
and~\eqref{eq:trunc-conv}:
\begin{align}
  \ip{f * h}{g}
  &= \sum_{m=0}^{k-1} \Bigl(\sum_{\substack{i+j=m\\0\le i,j<k}}
      f_i\, h_j\Bigr) g_m \notag\\
  &= \sum_{i=0}^{k-1} f_i \sum_{j=0}^{k-1-i} h_j\, g_{i+j}
  \label{eq:adj-swap}\\
  &= \sum_{i=0}^{k-1} f_i \cdot \bigl(\corr(g,h)\bigr)_i
  = \ip{f}{\corr(g,h)}\,. \notag
\end{align}
The exchange of summation in~\eqref{eq:adj-swap} is justified by the
finiteness of all sums, and the upper limit $k{-}1{-}i$ on the $j$-sum
arises from the truncation $i+j < k$.
\end{proof}

\begin{remark}[Relation to the Toeplitz transpose]
\label{rem:toeplitz}
The linear map $\mu_h$ has the Toeplitz matrix representation
$M_{ij} = h_{i-j}$ (with $h_\ell = 0$ for $\ell < 0$ or $\ell \ge k$).
Lemma~\ref{lem:adj-conv} asserts that the adjoint $\mu_h^*$ is the
transposed Toeplitz matrix $M^T_{ij} = h_{j-i}$, whose action on a vector
is precisely cross-correlation. This is the finite, truncated analogue of
the classical fact that the adjoint of a convolution operator on
$L^2(\mathbb{R})$ is convolution with the time-reversed kernel.
\end{remark}

% ======================================================================
\section{Algorithm Description}\label{sec:algorithm}
% ======================================================================

\subsection{Overview}

My algorithm proceeds in three phases for each quadrature point $v_q$:
\begin{enumerate}
  \item \textbf{Compute $a_j(v_q)$:} For each player $j$, compute
    $a_j = v_q^{S_j}$ and $b_j = 1 - a_j$.
  \item \textbf{Evaluate inner products:} Compute
    $I_i(v_q) = \sum_{r=1}^{k} \pi_r \cdot [x^{n-r}]\, Q_i(x; v_q)$
    for all players $i$ (Equation~\ref{eq:inner-product}).
  \item \textbf{Accumulate:} Add the weighted contribution
    $w_q \cdot S_i \cdot a_i \cdot v_q^{-1} \cdot I_i(v_q)$ to
    $E_i$.
\end{enumerate}

Phase~2 is the computational bottleneck. I provide three engines for this
phase, with automatic cost-model-driven dispatch.

\subsection{Engine 1: Linear Forward-Backward}\label{sec:linear}

The linear engine computes all $n$ inner products in $\bigO(nk)$ time using
a forward-backward decomposition.

\begin{algorithm}[H]
\caption{Linear Forward-Backward Engine}\label{alg:linear}
\DontPrintSemicolon
\SetAlgoLined
\KwIn{$n$ players, values $a_1, \ldots, a_n$, payout vector
  $\bm{\pi}$ with $k$ positions}
\KwOut{Inner products $I_1, \ldots, I_n$}
\BlankLine
\tcp{Forward pass: compute prefix products}
$g^{(0)}_m \leftarrow \pi_{n-m}$ for $m = 0, \ldots, k-1$\;
\For{$j = 1, \ldots, n$}{
  \For{$m = 0, \ldots, k-2$}{
    $g^{(j)}_m \leftarrow a_j \cdot g^{(j-1)}_m + (1-a_j) \cdot g^{(j-1)}_{m+1}$\;
  }
  $g^{(j)}_{k-1} \leftarrow a_j \cdot g^{(j-1)}_{k-1}$\;
}
\BlankLine
\tcp{Backward pass: compute suffix products and extract equities}
$R_0 \leftarrow 1$;\quad $R_m \leftarrow 0$ for $m = 1, \ldots, k-1$\;
\For{$j = n, n-1, \ldots, 1$}{
  $I_j \leftarrow \sum_{m=0}^{k-1} g^{(j-1)}_m \cdot R_m$ \tcp*{Inner product}
  \For{$m = k-1, k-2, \ldots, 1$}{
    $R_m \leftarrow a_j \cdot R_m + (1-a_j) \cdot R_{m-1}$\;
  }
  $R_0 \leftarrow a_j \cdot R_0$\;
}
\Return{$I_1, \ldots, I_n$}
\end{algorithm}

\begin{lemma}[Correctness of the fused backward pass]\label{lem:fused}
In Algorithm~\ref{alg:linear}, the inner product on line~8 and the suffix
update on lines~9--11 can be fused into a single loop because the backward
iteration over $m$ reads $R_m$ for the dot product before overwriting it:
$R_{m-1}$ is still pristine when used in the update of $R_m$.
\end{lemma}
\begin{proof}
At iteration $m$ (decreasing), the update sets $R_m \leftarrow a_j R_m +
(1-a_j) R_{m-1}$. The value $R_{m-1}$ has not yet been updated at this
iteration (it is updated at $m-1$), and $R_m$ was already read for the dot
product accumulation $I_j \mathrel{+}= g^{(j-1)}_m \cdot R_m$. Thus the
fused loop computes both quantities correctly.
\end{proof}

\paragraph{Batched quadrature.}
To exploit SIMD parallelism, I interleave $B_Q = 8$ quadrature points
in the data layout. The array $a_{\text{batch}}[j \cdot B_Q + q_i]$ stores
the $q_i$-th quadrature value for player $j$, ensuring cache-friendly
sequential access in the inner loops. The inner loop over the $B_Q$
quadrature dimension auto-vectorizes to full-width FMA instructions: 4
NEON ports $\times$ 128-bit on Apple Silicon, or native 512-bit on AVX-512.

\paragraph{L2-aware checkpointing.}
When the working set $n \cdot k \cdot B_Q \cdot 8$ bytes exceeds the L2
cache, I employ a checkpointing scheme inspired by gradient checkpointing
in automatic differentiation. The forward pass saves checkpoints every $C$
players (where $C$ is sized so one segment fits in L2), then the backward
pass recomputes each segment from its checkpoint before extracting equities.
This trades $\bigO(\sqrt{n})$ recomputation for $\bigO(\sqrt{n} \cdot k)$
memory.

\subsection{Engine 2: FFT-Accelerated Subproduct Tree}\label{sec:tree}

For large $k$, the $\bigO(nk)$ linear engine becomes expensive. I replace
it with a tree-based approach whose cost is $\bigO(n \log^2 k)$.

\begin{definition}[Binary subproduct tree]
Given $n$ linear factors $P_j(x) = a_j + b_j x$ for $j = 1, \ldots, n$,
the binary subproduct tree has $L = \lceil \log_2 n \rceil + 1$ levels.
Level~0 contains the $n$ leaves $P_j$. Level~$\ell$ contains
$\lceil n / 2^\ell \rceil$ nodes, each being the product of its two
children (truncated to degree $k$):
\[
  P^{(\ell)}_i = P^{(\ell-1)}_{2i} \cdot P^{(\ell-1)}_{2i+1} \bmod x^k\,.
\]
\end{definition}

\begin{algorithm}[H]
\caption{Tree Engine: Build and Propagate}\label{alg:tree}
\DontPrintSemicolon
\SetAlgoLined
\KwIn{Leaf polynomials $P^{(0)}_1, \ldots, P^{(0)}_n$,
  payout vector $\bm{\pi}$, truncation degree $k$}
\KwOut{Leaf-level $g$-vectors $g^{(0)}_1, \ldots, g^{(0)}_n$}
\BlankLine
\tcp{Build phase (bottom-up): compute subproducts}
\For{$\ell = 1, \ldots, L-2$}{
  \For{each parent node $i$ at level $\ell$ \textbf{in parallel}}{
    $P^{(\ell)}_i \leftarrow P^{(\ell-1)}_{2i} \cdot P^{(\ell-1)}_{2i+1}
      \bmod x^k$
      \tcp*{FFT or schoolbook}
  }
}
\BlankLine
\tcp{Propagation phase (top-down): distribute payout vector}
$g^{(L-1)}_m \leftarrow \pi_{n-m}$\;
\For{$\ell = L-1, L-2, \ldots, 1$}{
  \For{each parent node $i$ at level $\ell$ \textbf{in parallel}}{
    $g^{(\ell-1)}_{2i} \leftarrow \corr\bigl(g^{(\ell)}_i,\;
      P^{(\ell-1)}_{2i+1}\bigr)$
      \tcp*{Left child via right sibling}
    $g^{(\ell-1)}_{2i+1} \leftarrow \corr\bigl(g^{(\ell)}_i,\;
      P^{(\ell-1)}_{2i}\bigr)$
      \tcp*{Right child via left sibling}
  }
}
\Return{$g^{(0)}_1[0], \ldots, g^{(0)}_n[0]$}
  \tcp*{$I_j = g^{(0)}_j[0]$}
\end{algorithm}

\subsection{Correctness of the Tree Engine via the Adjoint}\label{sec:tree-correct}

The correctness of Algorithm~\ref{alg:tree} follows cleanly from the
adjoint identity of Lemma~\ref{lem:adj-conv}. The key computational
task is: for each player $j$, compute
\begin{equation}\label{eq:target-ip}
  I_j = \ip{\tilde\pi}{Q_j}
  = \sum_{m=0}^{k-1} \tilde\pi_m \cdot [x^m]\, Q_j(x;v)\,,
\end{equation}
where $\tilde\pi_m = \pi_{n-m}$ is the reversed payout vector and
$Q_j = \prod_{i \ne j} P_i \bmod x^k$ is the leave-one-out product.
The challenge is to evaluate all $n$ such inner products simultaneously
without forming each $Q_j$ independently.

\paragraph{The build phase as a composition of convolutions.}
The build phase constructs the full product
$R = P_1 * P_2 * \cdots * P_n \bmod x^k$
via a binary tree of truncated convolutions. Each internal node at
level~$\ell$ computes
$P^{(\ell)}_i = \mu_{P^{(\ell-1)}_{2i+1}}\!\bigl(P^{(\ell-1)}_{2i}\bigr)$,
i.e., ``convolve the left child with the right child.'' Reading from root
to any leaf~$j$, the full product can be decomposed as
\begin{equation}\label{eq:product-decomp}
  R = Q_j * P_j \bmod x^k\,,
\end{equation}
and $Q_j$ is the product of all the \emph{sibling} polynomials encountered
on the path from leaf~$j$ to the root.

\paragraph{The propagation phase as the adjoint of the build.}
The propagation phase computes $\ip{\tilde\pi}{Q_j}$ for all $j$ by
applying the adjoint of the build phase's linear map, with $\tilde\pi$
as the ``seed.'' At each internal node $i$ at level $\ell$, the build
phase applied the bilinear map
$(L, R) \mapsto L * R \bmod x^k$.
For the purpose of computing $\ip{\tilde\pi}{Q_j}$, we hold the sibling
polynomial fixed and view this as a \emph{linear} map in the descendant
of~$j$. The adjoint of $f \mapsto \mu_h(f) = f * h \bmod x^k$ is
$g \mapsto \mu_h^*(g) = \corr(g, h)$ by Lemma~\ref{lem:adj-conv}.

This is exactly what Algorithm~\ref{alg:tree} implements: each child
receives the correlation of its parent's $g$-vector with its
\emph{sibling}'s subproduct polynomial.

\begin{theorem}[Correctness of the tree engine]\label{thm:tree-correct}
After Algorithm~\ref{alg:tree} completes, $g^{(0)}_j[0] = I_j$ for all
$j = 1, \ldots, n$.
\end{theorem}
\begin{proof}
We establish the following invariant by downward induction on $\ell$.

\textbf{Invariant.} Let $\mathrm{leaves}(i,\ell)$ denote the set of leaf
indices in the subtree rooted at node $i$ at level $\ell$, and let
$C_j^{(\ell)} \subseteq [n] \setminus \mathrm{leaves}(i,\ell)$ be the set
of all leaves \emph{outside} this subtree.  After the propagation phase
processes level $\ell$, the $g$-vector at node $i$ satisfies
\begin{equation}\label{eq:prop-invariant}
  g^{(\ell)}_i = \corr\!\Bigl(\tilde\pi,\;
    \prod_{j \in C_j^{(\ell)}} P_j \bmod x^k \Bigr).
\end{equation}
In particular, $\ip{f}{g^{(\ell)}_i} =
\ip{f * \prod_{j \in C_j^{(\ell)}} P_j}{\tilde\pi}$ for all
$f \in \mathbb{R}[x]_{<k}$, by iterating Lemma~\ref{lem:adj-conv}.

\textbf{Base case} ($\ell = L{-}1$, the root). The root's subtree
contains all $n$ players, so $C_j^{(L-1)} = \emptyset$ and the product
over the empty set is~$1$.  The initialization
$g^{(L-1)}_m = \tilde\pi_m$ is $\corr(\tilde\pi, 1) = \tilde\pi$,
satisfying~\eqref{eq:prop-invariant}.

\textbf{Inductive step.} Assume the invariant holds at level $\ell$
for node $i$.  Write $C = C_j^{(\ell)}$ and let $R_C = \prod_{j \in C}
P_j \bmod x^k$.  By the inductive hypothesis, $g^{(\ell)}_i =
\corr(\tilde\pi, R_C)$.  The left child $2i$ receives
\[
  g^{(\ell-1)}_{2i}
  = \corr\bigl(g^{(\ell)}_i,\; P^{(\ell-1)}_{2i+1}\bigr)
  = \corr\!\bigl(\corr(\tilde\pi, R_C),\; P^{(\ell-1)}_{2i+1}\bigr).
\]
We claim this equals $\corr(\tilde\pi,\; R_C \cdot P^{(\ell-1)}_{2i+1}
\bmod x^k)$.  To verify, take an arbitrary $f \in \mathbb{R}[x]_{<k}$
and compute:
\begin{align*}
  \ip{f}{g^{(\ell-1)}_{2i}}
  &= \ip{f}{\corr\bigl(\corr(\tilde\pi, R_C),\; P^{(\ell-1)}_{2i+1}\bigr)}
  \\
  &= \ip{f * P^{(\ell-1)}_{2i+1}}{\corr(\tilde\pi, R_C)}
    &&\text{(Lemma~\ref{lem:adj-conv})}
  \\
  &= \ip{(f * P^{(\ell-1)}_{2i+1}) * R_C}{\tilde\pi}
    &&\text{(Lemma~\ref{lem:adj-conv})}
  \\
  &= \ip{f * (P^{(\ell-1)}_{2i+1} * R_C)}{\tilde\pi}
    &&\text{(associativity of $*$ mod $x^k$)}
  \\
  &= \ip{f}{\corr(\tilde\pi,\; P^{(\ell-1)}_{2i+1} * R_C)}
    &&\text{(Lemma~\ref{lem:adj-conv})}\,.
\end{align*}
Since this holds for all $f$, we conclude
$g^{(\ell-1)}_{2i} = \corr(\tilde\pi,\; R_C \cdot P^{(\ell-1)}_{2i+1}
\bmod x^k)$.  But $C_j^{(\ell-1)}$ for the left child $2i$ is
$C \cup \mathrm{leaves}(2i{+}1, \ell{-}1)$, and
$P^{(\ell-1)}_{2i+1} = \prod_{j \in \mathrm{leaves}(2i+1,\ell-1)} P_j
\bmod x^k$.  Thus $R_C \cdot P^{(\ell-1)}_{2i+1} =
\prod_{j \in C_j^{(\ell-1)}} P_j \bmod x^k$, confirming the invariant
for the left child.  The right child is symmetric.

\textbf{Conclusion.} At level~$0$, leaf~$j$ has
$C_j^{(0)} = [n] \setminus \{j\}$, so
$g^{(0)}_j = \corr(\tilde\pi, Q_j)$ where $Q_j = \prod_{i \ne j} P_i
\bmod x^k$.  Evaluating at index~$0$:
\[
  g^{(0)}_j[0]
  = \bigl(\corr(\tilde\pi, Q_j)\bigr)_0
  = \sum_{m=0}^{k-1} \tilde\pi_m \cdot (Q_j)_m
  = \ip{\tilde\pi}{Q_j}
  = I_j\,.\qedhere
\]
\end{proof}

\begin{remark}[Adjoint chain rule and complexity parity]
\label{rem:adjoint-chain}
The propagation phase is an instance of \emph{reverse-mode automatic
differentiation} applied to the build phase.  More precisely, define
the linear map $\Phi_j \colon \mathbb{R}[x]_{<k} \to
\mathbb{R}[x]_{<k}$ at each build node as $\Phi_j = \mu_{P_{\mathrm{sib}(j)}}$
(convolution with the sibling).  The build computes the composition
$\Phi_L \circ \cdots \circ \Phi_1$, and the propagation computes
$\Phi_1^* \circ \cdots \circ \Phi_L^*$ applied to $\tilde\pi$---the
adjoint of the composition, in reversed order, by the chain rule
$(\Phi_L \circ \cdots \circ \Phi_1)^* = \Phi_1^* \circ \cdots \circ
\Phi_L^*$.

A standard result in reverse-mode AD states that the adjoint of a linear
program has the same arithmetic complexity as the forward program
(Griewank and Walther, \emph{Evaluating Derivatives}, Theorem~3.1).
This immediately implies that the propagation cost equals the build
cost---a fact we use in Theorem~\ref{thm:tree-complexity}.
\end{remark}

\paragraph{Per-level FFT vs.\ schoolbook decisions.}
At each tree level, I compare the estimated cost of FFT-based and
schoolbook-based polynomial multiplication using offline-calibrated data.
Let $d_{\text{eff}}$ be the effective degree of the child polynomials at
level $\ell$ (which equals $\min(2^{\ell-1}, k) - 1$ at saturated levels
and half that at below-saturation levels). The decision rule is:
\begin{equation}\label{eq:fft-decision}
  \text{use FFT} \iff
    T_{\text{FFT}}(S^*) + T_{\text{overhead}} + (m+1)^2 \cdot \tau_{\text{FMA}}
    < (d_{\text{eff}}+1)^2 \cdot \tau_{\text{FMA}}\,,
\end{equation}
where $T_{\text{FFT}}(S^*)$ is the calibrated FFT time for the optimal
smooth size $S^*$, $T_{\text{overhead}}$ is the measured per-call FFT
overhead (plan lookup, buffer copies), $m$ is the wrap correction count, and
$\tau_{\text{FMA}}$ is the measured FMA latency.

\paragraph{Cyclic FFT with wrap correction.}
Rather than padding to the next power of two, I search all 7-smooth numbers
(of the form $2^a \cdot 3^b \cdot 5^c \cdot 7^d$) for the FFT size $S$
minimizing the total cost:
\begin{equation}\label{eq:wrap-cost}
  \text{cost}(S) = T_{\text{FFT}}(S) + (m+1)^2 \cdot \tau_{\text{FMA}}\,,
  \quad m = \max(L_{\text{conv}} - S, 0)\,,
\end{equation}
where $L_{\text{conv}}$ is the required convolution length. When $S <
L_{\text{conv}}$, the cyclic convolution of size $S$ wraps $m$ high-degree
terms to low-degree positions. These are corrected by a schoolbook product
of the top $(m+1)$ coefficients of each input.

\paragraph{Paired cached correlate.}
During propagation, both children at each node require correlating the same
parent $g$-vector with their respective sibling polynomials. I share the
forward FFT of $g$ across both correlates, saving one FFT per parent node.
Furthermore, when the build and correlate phases use the same FFT size
(determined by a joint cost optimization), I cache the FFT of each child
polynomial during the build phase and reuse it via conjugation during
propagation.

\paragraph{Ragged tree.}
I track the number of real (non-padding) nodes $n_{\text{real}}[\ell]$ at
each level. Lone children (when $n_{\text{real}}$ is odd at a level) are
copied up without multiplication. Padding nodes are never processed. This
saves approximately $1 - n/2^{\lceil \log_2 n \rceil}$ of the work at
non-power-of-two $n$.

\subsection{Engine 3: Hybrid Block-Divide}\label{sec:hybrid}

The hybrid engine partitions the $n$ players into $\lceil n/B \rceil$
blocks of size $B$ and combines a direct $\bigO(B^2)$ within-block
computation with the tree engine for inter-block structure.

\begin{algorithm}[H]
\caption{Hybrid Block-Divide Engine}\label{alg:hybrid}
\DontPrintSemicolon
\SetAlgoLined
\KwIn{$n$ players with values $a_1, \ldots, a_n$, block size $B$,
  payout vector $\bm{\pi}$, degree $k$}
\KwOut{Inner products $I_1, \ldots, I_n$}
\BlankLine
\tcp{Step 1: Build block products}
\For{each block $b = 0, 1, \ldots, \lceil n/B \rceil - 1$}{
  $P_b(x) \leftarrow \prod_{j \in \text{block } b} (a_j + b_j x)$
    \tcp*{Sequential, $\bigO(B^2)$}
}
\BlankLine
\tcp{Step 2: Inter-block tree (Alg.~\ref{alg:tree}) with leaves $P_b$}
Build subproduct tree from block products $P_0, \ldots, P_{\lceil n/B \rceil - 1}$\;
Propagate $\bm{\pi}$ to obtain block-level $g$-vectors $g_b$\;
\BlankLine
\tcp{Step 3: Within-block divide}
\For{each block $b$}{
  \For{each player $j$ in block $b$}{
    Compute $Q_j^{(b)} \leftarrow P_b(x) / (a_j + b_j x)$
      via stable bidirectional synthetic division\;
    $I_j \leftarrow \sum_{m=0}^{\min(k,B)-1} g_{b,m} \cdot Q_{j,m}^{(b)}$\;
  }
}
\Return{$I_1, \ldots, I_n$}
\end{algorithm}

\begin{remark}[Bidirectional division stability]
Synthetic division of $P_b(x)$ by $(a_j + b_j x)$ is numerically unstable
when $a_j < 0.5$ (forward direction amplifies errors). I use bidirectional
division: forward when $a_j > 0.5$, backward when $a_j \leq 0.5$,
guaranteeing that the division coefficient satisfies $|c| < 1$ in all cases.
Players are sorted by stack size so that within each block, most players
use the same direction, improving branch prediction.
\end{remark}

\paragraph{Cost-model block size selection.}
The block size $B$ is selected by a cost model that evaluates candidates
$B \in \{8, 16, 24, 32, 48, 64\}$ and picks the one minimizing:
\begin{equation}\label{eq:hybrid-cost}
  T_{\text{hybrid}}(B) = T_{\text{block}}(B) + T_{\text{tree}}(n/B, k)\,,
\end{equation}
where $T_{\text{block}}$ accounts for the $\bigO(n B)$ block build and
divide cost, and $T_{\text{tree}}$ is estimated using the actual per-level
planning decisions (FFT sizes, schoolbook/FFT dispatch, wrap corrections)
from a lightweight trial context creation. Typical selections: $B = 16$ on
M3~Max, $B = 32$ on Zen~4.

\subsection{Engine Dispatch}\label{sec:dispatch}

\begin{algorithm}[H]
\caption{Cost-Based Engine Dispatch}\label{alg:dispatch}
\DontPrintSemicolon
\SetAlgoLined
\KwIn{$n, k$, calibration constants}
\KwOut{Engine choice and parameters}
\BlankLine
$c_{\text{ckpt}} \leftarrow
  \begin{cases} 1.0 & \text{if } n \cdot k \cdot 64 \leq L_2 \\
                1.3 & \text{otherwise}
  \end{cases}$\;
$T_{\text{linear}} \leftarrow 2nk \cdot \tau_{\text{poly}} \cdot
  c_{\text{ckpt}}$\;
$B^* \leftarrow \text{select\_best\_B}(n, k)$\;
$T_{\text{hybrid}} \leftarrow T_{\text{block}}(B^*) +
  T_{\text{tree}}(n/B^*, k)$\;
\lIf{$T_{\text{hybrid}} < T_{\text{linear}}$}{
  \Return{Hybrid with $B = B^*$}
}
\lElse{
  \Return{Linear batched}
}
\end{algorithm}

The dispatch is fully derived from calibration constants
($\tau_{\text{FMA}}$, $\tau_{\text{poly}}$, $T_{\text{FFT}}(\cdot)$,
$L_2$) measured on the target hardware. No manual crossover thresholds
are needed.

% ======================================================================
\section{Theoretical Analysis}\label{sec:theory}
% ======================================================================

\subsection{Time Complexity}

We analyze complexity by establishing recurrences for the build and
propagation phases and solving them via the Master theorem, rather than
summing per-level costs directly.

\begin{theorem}[Complexity of the tree engine]\label{thm:tree-complexity}
The tree engine (Algorithm~\ref{alg:tree}) computes all $n$ inner products
in $\bigO(n \log^2 k)$ arithmetic operations per quadrature point.
\end{theorem}
\begin{proof}
Let $T_b(n)$ denote the total arithmetic cost of the build phase for $n$
leaves with truncation degree $k$.  At the root, two subtrees of
$\lfloor n/2 \rfloor$ and $\lceil n/2 \rceil$ leaves are combined by a
single truncated polynomial multiplication.  Each child's subproduct has
degree $\min(\lfloor n/2 \rfloor, k)$, so the multiplication costs
$\bigO(\min(n, k) \cdot \log \min(n, k))$ via FFT.  This gives the
recurrence
\begin{equation}\label{eq:build-recurrence}
  T_b(n) = 2\, T_b(n/2)
  + \bigO\!\bigl(\min(n, k) \cdot \log \min(n, k)\bigr),
  \qquad T_b(1) = 0.
\end{equation}

\textbf{Case 1: $n \le k$ (below saturation).}
The recurrence becomes $T_b(n) = 2\,T_b(n/2) + \bigO(n \log n)$.
The merge cost $f(n) = n \log n$ satisfies
$f(n) = \Theta(n^{\log_2 2} \log n) = \Theta(n \log n)$, which is
Case~2 of the Master theorem (with $p = 1$), yielding
\begin{equation}\label{eq:below-sat}
  T_b(n) = \bigO(n \log^2 n)\,.
\end{equation}

\textbf{Case 2: $n > k$ (above saturation).}
All subproducts are capped at degree $k$, so the merge cost is
$\bigO(k \log k)$ regardless of $n$.  The recurrence becomes
$T_b(n) = 2\, T_b(n/2) + \bigO(k \log k)$ with base case
$T_b(k) = \bigO(k \log^2 k)$ from Case~1.
Unrolling for $s = \log_2(n/k)$ levels above saturation:
\begin{align}
  T_b(n) &= \frac{n}{k}\, T_b(k)
    + \bigO(k \log k) \sum_{i=0}^{s-1} 2^i \notag\\
  &= \frac{n}{k} \cdot \bigO(k \log^2 k)
    + \bigO\!\bigl(\tfrac{n}{k} \cdot k \log k\bigr) \notag\\
  &= \bigO(n \log^2 k) + \bigO(n \log k)
  = \bigO(n \log^2 k)\,. \label{eq:above-sat}
\end{align}

In both cases, $T_b(n) = \bigO(n \log^2 k)$.  By
Remark~\ref{rem:adjoint-chain}, the propagation phase has the same
arithmetic complexity as the build phase, so the total per-quadrature-point
cost is $T_b(n) + T_p(n) = \bigO(n \log^2 k)$.  Over $Q$ quadrature
points, the total is $\bigO(Qn \log^2 k)$.
\end{proof}

\begin{theorem}[Complexity of the hybrid engine]\label{thm:hybrid-complexity}
The hybrid engine (Algorithm~\ref{alg:hybrid}) with block size $B$
computes all $n$ inner products in
$\bigO(nB + (n/B) \log^2 k)$ operations per quadrature point.
\end{theorem}
\begin{proof}
The block build step produces $n/B$ block polynomials, each of
degree~$B$, at a cost of $\bigO(B^2)$ per block (schoolbook multiplication
of $B$ linear factors), totaling $\bigO(nB)$.
The inter-block tree over $n/B$ leaves with truncation degree $k$ costs
$\bigO((n/B) \log^2 k)$ by Theorem~\ref{thm:tree-complexity}.
The within-block divide performs, for each of $n$ players, a synthetic
division of a degree-$B$ polynomial and a dot product of length $B$,
costing $\bigO(B)$ per player and $\bigO(nB)$ in total.
Summing: $\bigO(nB + (n/B) \log^2 k)$.
\end{proof}

\begin{corollary}[Optimal block size]\label{cor:opt-block}
Setting $B = \Theta(\log k)$ balances the two terms in
Theorem~\ref{thm:hybrid-complexity}, yielding $\bigO(n \log^2 k)$ total
cost--the same asymptotic complexity as the pure tree engine.  In practice,
hardware constraints (SIMD width, cache line size, schoolbook-FFT crossover)
determine the optimal $B$ via the cost model of
Equation~\eqref{eq:hybrid-cost}.
\end{corollary}

\begin{proposition}[Complexity of the linear engine]\label{prop:linear-complexity}
The linear forward-backward engine (Algorithm~\ref{alg:linear}) computes
all $n$ inner products in $\bigO(nk)$ operations per quadrature point.
\end{proposition}

This is immediate from inspection: the forward pass performs $n$ iterations
of $\bigO(k)$ work each, and the backward pass is identical.

\begin{remark}[Dispatch optimality]
The linear engine is asymptotically optimal for $k = \bigO(1)$ (fixed
payout depth). The hybrid/tree engines are optimal for $k = \Theta(n)$
(all players receive payouts). The cost-model dispatch automatically
selects the faster engine for each $(n, k)$ pair.
\end{remark}

\subsection{Numerical Accuracy}

The exponential convergence of the trapezoidal rule on analytic integrands
is a classical result; we state it here without proof and refer the reader
to Trefethen and Weideman~(2014) for a comprehensive treatment.

\begin{theorem}[Quadrature convergence; Trefethen--Weideman 2014]
\label{thm:quad}
Let $f$ be analytic in a strip $|\operatorname{Im}(y)| < d$ around the
real axis and satisfy $\int_{-\infty}^{\infty} |f(y)|\, dy < \infty$.
The trapezoidal rule with step size $h$ and $Q$ points satisfies
\[
  \biggl|\int_{-\infty}^{\infty} f(y)\, dy
  - h \sum_{q} f(y_q)\biggr|
  = \bigO\!\bigl(e^{-2\pi d/h}\bigr).
\]
Under the change of variables $v = \Phi(y)$, the ICM integrand is analytic
in a strip whose width depends on the chip stacks.  With $Q = 256$ and
typical strip widths, the quadrature error is well below double-precision
machine epsilon ($\approx 10^{-16}$), consistent with observed relative
errors $< 5 \times 10^{-12}$.
\end{theorem}

\begin{lemma}[Numerical stability of the hybrid divide]\label{lem:divide}
The bidirectional synthetic division in Algorithm~\ref{alg:hybrid} has
per-step error amplification factor strictly less than~$1$ for all
$a_j \in (0, 1)$.
\end{lemma}
\begin{proof}
Forward division computes $Q_m = (P_m - (1{-}a_j) Q_{m-1}) / a_j$.
A perturbation $\delta$ in $Q_{m-1}$ propagates as
$(1{-}a_j)/a_j \cdot \delta$.  When $a_j > 1/2$, the amplification
factor $(1{-}a_j)/a_j < 1$, so errors contract geometrically.

Backward division computes $Q_m = (P_{m+1} - a_j Q_{m+1}) / (1{-}a_j)$,
with amplification factor $a_j/(1{-}a_j)$.  When $a_j \le 1/2$, this
factor is at most~$1$.

Choosing forward division when $a_j > 1/2$ and backward division when
$a_j \le 1/2$ ensures the amplification factor is strictly less than~$1$
at every step, so accumulated rounding errors decay exponentially along
the division chain.
\end{proof}

% ======================================================================
\section{Experimental Results}\label{sec:experiments}
% ======================================================================

\subsection{Experimental Setup}

I evaluate my implementation on two platforms:
\begin{itemize}
  \item \textbf{Apple M3~Max:} 12 performance + 4 efficiency cores,
    ARM64 with NEON SIMD (2 FP64/vector), 4 FMA ports per core,
    32\,MB shared L2 cache, 400\,GB/s unified memory bandwidth.
    FFT backend: FFTW3 with PATIENT wisdom + Apple vDSP interleaved DFT
    at supported sizes.
  \item \textbf{AMD Ryzen 9 7950X (Zen 4):} 16 physical cores,
    x86-64 with AVX-512 (8 FP64/vector), 2 FMA units per core,
    1\,MB per-core L2 cache, ${\sim}60$\,GB/s DDR5 memory bandwidth.
    FFT backend: FFTW3 with PATIENT wisdom + Intel MKL via \texttt{dlopen}
    dual dispatch.
\end{itemize}

All experiments use $Q = 256$ quadrature points, uniform chip stacks and report the median of 5
repetitions. Correctness is verified against closed-form $V_1$ and $V_2$
equities and cross-checked between engines, with relative errors below
$5 \times 10^{-12}$.

\subsection{Performance Results}

Table~\ref{tab:serial} reports single-threaded performance. The engine
dispatch automatically selects the linear engine for small $k$ and the
hybrid engine for large $k$.

\begin{table}[t]
\centering
\caption{Single-threaded execution time (ms) with $Q = 256$.
  Engine selection shown as L (linear) or H (hybrid).}
\label{tab:serial}
\begin{tabular}{@{}l rrrrrr rrrrrr@{}}
\toprule
& \multicolumn{6}{c}{\textbf{Apple M3 Max}} &
  \multicolumn{6}{c}{\textbf{AMD Zen 4}} \\
\cmidrule(lr){2-7} \cmidrule(lr){8-13}
$n$ & $k{=}10$ & $k{=}50$ & $k{=}100$ & $k{=}n/4$ & $k{=}n/2$ & $k{=}n$
    & $k{=}10$ & $k{=}50$ & $k{=}100$ & $k{=}n/4$ & $k{=}n/2$ & $k{=}n$ \\
\midrule
1024  & 3 & 11 & 16 & 21 & 24 & 27 & 5 & 9 & 15 & 19 & 22 & 24 \\
2048  & 4 & 12 & 23 & 51 & 58 & 63 & 3 & 7 & 15 & 47 & 51 & 55 \\
4096  & 7 & 24 & 46 & 121 & 137 & 147 & 7 & 14 & 28 & 109 & 119 & 130 \\
8192  & 14 & 50 & 115 & 287 & 318 & 350 & 15 & 29 & 56 & 251 & 287 & 296 \\
16384 & 28 & 122 & 230 & 660 & 709 & 752 & 29 & 65 & 114 & 604 & 640 & 690 \\
32768 & 55 & 240 & 457 & 1511 & 1710 & 1808 & 71 & 132 & 219 & 1360 & 1517 & 1597 \\
65536 & 135 & 460 & 937 & 3480 & 4017 & 4392 & 142 & 259 & 510 & 3198 & 3473 & 3861 \\
\bottomrule
\end{tabular}
\end{table}

\begin{table}[t]
\centering
\caption{16-thread parallel execution time (ms) with $Q = 256$.}
\label{tab:parallel}
\begin{tabular}{@{}l rrrrrr rrrrrr@{}}
\toprule
& \multicolumn{6}{c}{\textbf{Apple M3 Max (16 threads)}} &
  \multicolumn{6}{c}{\textbf{AMD Zen 4 (16 threads)}} \\
\cmidrule(lr){2-7} \cmidrule(lr){8-13}
$n$ & $k{=}10$ & $k{=}50$ & $k{=}100$ & $k{=}n/4$ & $k{=}n/2$ & $k{=}n$
    & $k{=}10$ & $k{=}50$ & $k{=}100$ & $k{=}n/4$ & $k{=}n/2$ & $k{=}n$ \\
\midrule
1024  & $<1$ & 1 & 2 & 2 & 2 & 3 & $<1$ & 1 & 1 & 1 & 2 & 2 \\
4096  & 2 & 6 & 7 & 12 & 14 & 15 & 2 & 4 & 4 & 8 & 9 & 9 \\
8192  & 4 & 12 & 14 & 28 & 33 & 37 & 4 & 8 & 8 & 18 & 21 & 24 \\
16384 & 9 & 23 & 28 & 68 & 81 & 90 & 8 & 15 & 17 & 46 & 59 & 67 \\
65536 & 39 & 102 & 114 & 437 & 501 & 594 & 36 & 122 & 197 & 384 & 423 & 551 \\
\bottomrule
\end{tabular}
\end{table}

\paragraph{Parallel speedup.}
At $n = 8192$, $k = n$: M3~Max achieves 9.5$\times$ speedup
(350$\to$37\,ms) on 16 threads; Zen~4 achieves 12.3$\times$
(296$\to$24\,ms) on 16 threads. The higher speedup on Zen~4
reflects its homogeneous core topology (16 identical P-cores)
versus M3~Max's heterogeneous 12P+4E configuration.

\paragraph{Platform comparison.}
Zen~4 is ${\sim}2\times$ faster at small $k$ due to AVX-512's
8-wide FP64 FMA versus NEON's 2-wide, and 12--16\% faster at
large $k$ due to superior FFTW PATIENT wisdom on x86.

\subsection{Platform-Specific Optimizations}

\paragraph{Apple M3 Max: vDSP FFT dispatch.}
Apple's \texttt{vDSP\_DFT\_Interleaved} API provides 10--18\% faster FFTs
at 33 supported sizes ($f \times 2^g$ where $f \in \{1,3,5,15\}$,
$g \geq 4$). Crucially, this API uses the same interleaved complex memory
layout as FFTW, requiring only a 2-scalar fixup of the DC/Nyquist bin
packing per transform. The forward $\times 2$ scaling factor is absorbed
into the pointwise multiply (both operands carry the same factor), with a
single $\times 0.25$ correction on the inverse output--eliminating 2 of 3
per-round-trip scaling passes.

\paragraph{Apple M3 Max: AMX coprocessor.}
For schoolbook polynomial multiplication at degree $\geq 160$, I dispatch
to Apple's undocumented AMX coprocessor, which computes $8 \times 8$ FP64
outer products in ${\sim}2$\,ns per tile. Anti-diagonal extraction
(converting outer-product tiles to convolution coefficients) uses in-register
VECFP accumulation, which is $4.7\times$ faster than the STZ+scalar
fallback. The crossover at degree~160 is dictated by the per-column
extraction overhead (${\sim}69$\,ns) dominating at small sizes.

\paragraph{AMD Zen 4: FFTW/MKL dual dispatch.}
Both FFTW and Intel MKL are loaded via \texttt{dlopen} at initialization.
Offline calibration benchmarks the full FFT pipeline with both libraries at
each of 749 smooth sizes, recording the winner in a \texttt{calib\_lib[]}
array. At runtime, each FFT plan is created using the calibration-selected
library. MKL wins at 181/749 sizes (mostly small composites 14--64 and
131072), while FFTW dominates at 568/749 sizes with PATIENT wisdom.

\paragraph{AMD Zen 4: L2-aware checkpointing.}
Zen~4's 1\,MB per-core L2 (versus M3~Max's 32\,MB cluster L2) makes
checkpointing critical for the batched linear engine. The checkpoint
interval is sized so that one segment's working set
($C \times k \times B_Q \times 8$ bytes) fits in L2, ensuring the
recomputation pass runs at L2 bandwidth (${\sim}1$\,TB/s) rather than
DRAM bandwidth (${\sim}60$\,GB/s).

\subsection{Calibration Constants}

Table~\ref{tab:constants} summarizes the device-specific calibration constants measured
via offline profiling.

\begin{table}[t]
\centering
\caption{Device-specific calibration constants.}
\label{tab:constants}
\begin{tabular}{@{}lrrl@{}}
\toprule
\textbf{Constant} & \textbf{M3 Max} & \textbf{Zen 4} & \textbf{Description} \\
\midrule
$\tau_{\text{FMA}}$ (ns) & 0.25 & 0.13 & Scalar FMA latency \\
$\tau_{\text{poly}}$ (ns) & 0.13 & 0.135 & Memory-bound polymul FMA \\
$T_{\text{overhead}}$ (ns) & 40 & 50 & Per-call FFT overhead \\
Paired cache ratio & 1.03 & 1.08 & Paired correlate / pipeline \\
Indep.\ pair ratio & 1.25 & 1.35 & Independent pair / pipeline \\
$L_2$ cache (bytes) & 32\,MB & 1\,MB & Per-core L2 size \\
Typical $B$ & 16 & 32 & Cost-model-selected block size \\
FFT backend & FFTW+vDSP & FFTW+MKL & Library dispatch \\
\bottomrule
\end{tabular}
\end{table}

% ======================================================================
\section{GPU Acceleration: NVIDIA H200}\label{sec:gpu}
% ======================================================================

I outline the design for porting the tree-based engine to the NVIDIA H200
GPU (16896 CUDA cores, 4.8\,TB/s HBM3e bandwidth, ${\sim}34$\,TFLOPS FP64).
The linear engine is inherently sequential and does not benefit from GPU
parallelism; only the tree-based engines are applicable.

\subsection{Architecture Mapping}

The tree's level-by-level structure maps naturally to GPU execution: all
node operations at each level are independent and can execute as one batched
kernel. The build phase becomes $L-2$ batched cuFFT calls (each batching
$Q \times n_{\text{nodes}}[\ell]$ independent FFTs), and the propagation
phase is structured identically.

\paragraph{Batched cuFFT.}
At each tree level $\ell$, batch $Q \times \lceil n / 2^\ell \rceil$
independent FFTs into a single \texttt{cufftMakePlanMany} call.
For $n = 8192$, $Q = 256$, this yields up to $2^{20}$ simultaneous FFTs
at the bottom levels--sufficient to saturate the H200's bandwidth.

\paragraph{Schoolbook-FFT crossover.}
Prior work on GPU polynomial arithmetic suggests the
schoolbook$\to$FFT crossover on GPU occurs at degree ${\sim}4096$, far
higher than the ${\sim}32$ on CPU. This is because shared-memory schoolbook
kernels have excellent coalesced access patterns, while cuFFT has non-trivial
per-call overhead. This higher crossover fundamentally changes the hybrid
strategy: block sizes of $B = 64$--$256$ may be optimal (versus $B = 16$--$32$
on CPU), eliminating 6--8 bottom tree levels and their associated kernel
launch barriers.

\subsection{Optimization Strategies}

\paragraph{CUDA Graphs.}
The entire tree (build + propagation) can be captured as a CUDA Graph and
replayed with ${\sim}2.5\,\mu$s total launch overhead (versus
${\sim}3\,\mu$s per individual launch $\times$ 45 launches $\approx
135\,\mu$s). The tree structure is fixed per $(n, k)$, so the graph is
reused across all $Q$ quadrature points.

\paragraph{cuFFT LTO callbacks.}
Fusing the pointwise multiply into the inverse FFT's load callback
eliminates one full memory pass per tree node, reducing the
FFT-multiply-IFFT pipeline by ${\sim}20\%$ of memory traffic.

\paragraph{Memory layout.}
Polynomial data is organized as
$\texttt{poly[level][batch][coeff]}$ where
$\texttt{batch} = q \times n_{\text{nodes}}[\ell] + \text{node}$, ensuring
coalesced access for cuFFT. Memory footprint for $n = 8192$, $k = 8192$,
$Q = 256$: ${\sim}3.5$\,GB, fitting comfortably in the H200's 141\,GB HBM.

\subsection{Expected Performance}

\textbf{[Placeholder: benchmark results pending implementation.]}
Based on the H200's 4.8\,TB/s bandwidth and the algorithm's memory-access
pattern (each coefficient read/written ${\sim}3 \times L$ times across the
tree), I estimate sub-millisecond latency for $n = 8192$, $k = 8192$, a
${\sim}300\times$ speedup over the single-threaded CPU implementation.

% ======================================================================
\section{Conclusion}\label{sec:conclusion}
% ======================================================================

I have presented a high-performance algorithm for computing ICM tournament
equities using generating-function quadrature and FFT-accelerated subproduct
trees. The algorithm achieves $\bigO(Qn\log^2 k)$ time complexity, provides
double-precision accuracy via $Q = 256$ quadrature points, and
automatically dispatches between a linear engine (optimal for small $k$)
and a hybrid block-tree engine (optimal for large $k$) using a cost model
calibrated to the target hardware.

My implementation demonstrates the importance of hardware-aware algorithm
engineering: the same algorithmic framework achieves competitive performance
on two fundamentally different microarchitectures (ARM64 with vDSP/AMX
versus x86-64 with AVX-512/MKL) by parameterizing all dispatch decisions
through offline-calibrated constants. The key enablers are: (1)~per-size FFT
calibration that avoids the assumption that power-of-two sizes are always
fastest, (2)~cyclic convolution with schoolbook wrap correction that
exploits faster non-power-of-two FFT sizes, and (3)~paired cached
correlates that share both forward FFT computation and cached frequency-domain
polynomial data across sibling nodes.

Future work includes completing the CUDA port for the NVIDIA H200
(Section~\ref{sec:gpu}), extending the algorithm to non-uniform chip stack
distributions with adaptive quadrature, and exploring applications to
real-time tournament decision support where subset equity queries
(computing equities for a small number of ``hero'' players) can exploit the
active-player masking already present in my implementation.


\end{document}
